\cotefile{157-4}
\scanpage{303}\mspage{27}
D'autre part, je \unsure{soupçonne} que $\mathrm{loc}\,6\,a$ n'est
guère vérifié que si \deleted{\unsure{on suppose déjà}} $W=W^{\unsure{\leftrightarrows|\rightleftarrows}}$
\deleted{\illegible{1}} (cas qui est le plus commun bien sûr\unsure{\textsuperscript{2}}),
\inserted{auquel cas \deleted{\unsure{est}} $\mathrm{loc}\,6\,a$ \uline{équivaut} à $\mathrm{loc}\,6$}.
Il n'y a sans doute pas lieu de l'inclure
dans $\mathrm{loc}\,6$, pas plus que $W=W^{\leftrightarrows}$ dans $\mathrm{loc}\,1$.

\begin{center}\rule{0.15\linewidth}{0.4pt}\end{center}

4. \uline{Axiome} $\mathrm{L}\,7$ \uline{des $W$-fibrations}.

4.1. Rappelons la notion de \uline{$W$-fibration faible}
($\overline{\underline{\mathrm{XII}}}$ p. 51, 52) : la version la plus forte des \unsure{8}
conditions équivalentes pour que $f:X\mathrel{\text{---}}S$
soit une fibration faible est celle-ci
\[
\begin{aligned}
&\forall\, S''\longrightarrow S' \ \text{dans }\mathrm{Cat}_{/S},\quad
   \text{avec } S',S''\ \text{\uline{contractiles}},\\
&\text{le morphisme correspondant }\ X''=X\times_S S''\longrightarrow X'=X\times_S S'\\
&\text{est dans } W.
\end{aligned}
\tag*{\normalfont($\deleted{\unsure{4}}\mathrm{F}f$)}
\]
Quand $W$ est \unsure{fortement} saturé, alors ceci
implique la condition, en apparence plus forte :
\[
\begin{aligned}
&\forall\, S''\longrightarrow S'\ \text{dans }\mathrm{Cat}_{/S},\quad \text{\unsure{tels|tel} que }S''\mathrel{\text{---}}S'\\
&\text{soit un homotopisme},\quad X''\mathrel{\text{---}}X'\ \text{est dans } W.
\end{aligned}
\tag*{($\mathrm{F}f'$)}
\]
Je désigne par $\mathrm{Fib}f_W$ l'ensemble des
fibrations faibles. \deleted{\unsure{Les incl.}} Je résume dans
un diagramme d'inclusions les principales
classes de flèches de $\mathrm{Cat}$ que nous avons
rencontrées, s'apparentant à des classes
de ``fibrations'' (pour $W$, dans un sens
\unsure{convenable|honorable}) :
% archive sheet number: 302
\scanpage{304}\mspage{28}
\diagram{\begin{tikzcd}[ampersand replacement=\&,column sep=small,row sep=small]
\& \& \mathrm{Fib}_W \arrow[d] \& \& \\
\mathrm{Perf}_{\infty}\arrow[r] \& \mathrm{Perf}_W\arrow[r,"\deleted{\unsure{\leftarrow}}" description] \& \mathrm{Fib}f_W \&
\underbrace{\mathrm{Catfib}^{W}_{\deleted{\unsure{W}}}}_{\mathcal F_W}\arrow[l] \&
\underbrace{\mathrm{Catfib}^{\unsure{\infty|W}}_{\deleted{\unsure{W}}}}_{\mathcal F_\infty}\arrow[l,"\unsure{\Rightarrow}"']\;,
\end{tikzcd}}

où $\mathrm{Fib}_W$ désigne les $W$-fibrations,
($f:\deleted{\unsure{X}}S'\longrightarrow\deleted{\unsure{X}}S$ est dans $\mathrm{Fib}_W$ sss le foncteur
chgt de base $X\longmapsto X'=X\times_S S'$ de $\mathrm{Cat}_{/S}$ dans
$\mathrm{Cat}_{/S'}$ \deleted{\unsure{transforme les flèches}} \inserted{est compatible} avec les localiseurs
$W_S$ et $W_{S'}$), $\mathrm{Catfib}$ les Cat-fibrations,
$\mathrm{Catfib}^{W}_{\deleted{\unsure{W}}}$ \inserted{(resp. $\mathrm{Catfib}^{\infty}_{\deleted{\unsure{W}}}$)}
% slanted
\margin{NB\\ $\mathrm{Catfib}^{W}=\mathrm{Catfib}\cap\mathrm{Fib}f_W$\\
$\mathrm{Catfib}^{\infty}_{\unsure{\bullet|\circ}}=\mathrm{Catfib}^{\unsure{\circ}}\cap\mathrm{Fib}f_{W_\infty}$}
(les Cat-fibrations telles que les
foncteurs chgt. de base \deleted{\unsure{des fibres}} \deleted{\unsure{soient}} $\in W$
\inserted{(resp. soient $\in W_\infty$, cf. déf. plus bas)}),
$\mathrm{Perf}_W$ les morphismes $W$-parfaits, i.e.
$W$-lisses et $W$-propres, enfin $\mathrm{Perf}_\infty$ les
morphismes $W_\infty$-parfaits, où $W_\infty$ est
le plus petit localiseur satisfaisant
$\mathrm{Loc}(1,2,3)$. Parmi les \deleted{\unsure{six|cinq}} espèces de
``fibrations'' envisagées, c'est cette dernière
seulement, \inserted{avec $\mathrm{Catfib}^{\infty}_{\deleted{\unsure{W}}}$} qui est indépendante de $W$.
À ce titre, elle me paraît jouer un
rôle similaire, \inserted{dual,} à celui des morphismes
de Dwyer, \unsure{?} et je présume qu'on peut leur
faire jouer un rôle similaire. En tous
% archive sheet number: 303
\scanpage{305}\mspage{29}
cas, on sait que tte flèche $f$ de $\mathrm{Cat}$ se factorise
en $pi$, avec $i\in W_\infty$ et $p\in\mathrm{Perf}_\infty$, donc
$i\in W$ pour tt $W$ satisfaisant $\mathrm{loc}(1,2,3)$. \inserted{On
sait aussi que $\mathrm{Catfib}^{\unsure{\infty|W}}_{\deleted{\unsure{W}}}$ satisfait l'énoncé de factorisation,
tout $f$ se factorise en $f=pi$, $i\in W_\infty$, $p\in\mathrm{Catfib}^{\infty}_{\unsure{\circ|\bullet}}$.}
% slanted
\margin{donc la propriété de factorisation est valable pour \deleted{\unsure{toutes}} ces classes, sauf peut-être $\mathrm{Fib}_W$.}

Notons \unsure{aussi|ensuite} que $\mathrm{Perf}_\infty$ et $\mathrm{Perf}_W$, ainsi
que $\mathrm{Fib}_W$, sont stables non seulement par
chgt de base (c'est le cas des 5 classes
envisagées, qui de plus contiennent toutes
les iso), mais aussi par composition.
(De plus $\mathrm{Fib}_W$ est stable par facteurs
directs, et je présume qu'il en est de
\unsure{$=$|même} de $\deleted{\mathrm{Perf}_W}$ $\mathrm{Liss}_W$ et de $\mathrm{Prop}_W$, donc de
$\mathrm{Perf}_W$, en particulier de $\mathrm{Perf}_{W_\infty}$, mais
je n'ai pas pris la peine de le
vérifier). Par contre, a priori
$\mathrm{Fib}f_W$ n'est pas \inserted{en général} stable par composition,
et $\mathrm{Catfib}^{W}_{\deleted{\unsure{W}}}$, \deleted{\unsure{ne l'est pas, donc}} \inserted{$\mathrm{Catfib}^{\infty}_{\deleted{\unsure{W}}}$} n'ont pas de
raison non plus d'\deleted{\unsure{avoir}}\unsure{être} \deleted{cette propriété}
stables par composition.
% slanted
\margin{il faudrait l'expliciter par des exemples}

4.2. Finalement, il paraît plus ``éclairant'' d'introduire, au lieu du ``localiseur fondamental $W_\infty$
bien déterminé\unsure{''}, un localiseur fondamental
\[
W_0\subset\mathrm{Fl}(\mathrm{Cat}),
\]
qui jouera le rôle d'une sorte de ``localiseur de
référence'' pour formuler certaines propriétés de $W$.
On \deleted{\unsure{est}} supposera toujours que
\[
W\supset W_0.
\]
% archive sheet number: 304
\scanpage{306}\mspage{30}
Suivant les cas, $W_0$ pourra être $W_\infty$, ou quelque
autre localiseur fondamental défini comme le
``plus petit'' satisfaisant à des axiomes de stabilité
donnés. Ou ça pourra être $\unsure{\mathbb W|W}_\infty$ (l'équivalence
faible habituelle), quand on s'intéresse aux
localiseurs qui contiennent $\mathbb W_\infty$ (les seuls, apparemment, avec lesquels on ait travaillé jusqu'à présent en homotopie). Ça pourra être
aussi $W$, et les énoncés obtenus en spécialisant les résultats obtenus au cas où $\deleted{\unsure{W_0=W}}$
$W=W_0$ ne sont nullement tautologiques.

Le diagramme de la page 28 se remplace alors
par le diagramme \inserted{d'inclusions} suivant, plus éclairant
pour notre propos :
\diagram{\begin{tikzcd}[ampersand replacement=\&,column sep=small,row sep=small]
\& \& \mathrm{Fib}_W \arrow[d] \& \&\\
\& \mathrm{Perf}_W\arrow[r] \& \mathrm{Fib}f_W \& \mathcal F_W\arrow[l] \&\\
(4.2.1.) \& \mathrm{Perf}_{W_0}\arrow[u]\arrow[r,dash] \& \mathrm{Fib}f_{W_0}\arrow[u] \& \mathcal F_{W_0}\arrow[l]\arrow[u] \& \unsure{\ast|4}\\
\& \& \mathrm{Fib}_{W_0}\arrow[u] \& \& \text{,}
\end{tikzcd}}
où pour tout \deleted{\unsure{lo}}localiseur fondamental $\deleted{\unsure{\mathcal F}}W$, je
pose pour abréger
\[
\underline{\mathcal F}_W=\mathrm{Catfib}^{W}_{\deleted{\unsure{W}}}\,,
\]
ce qui définit donc également $\mathcal F_{W_0}$. La classe $\mathrm{Fib}_W$
% archive sheet number: 305
\scanpage{307}\mspage{31}
doit être regardée comme le point de mire de
notre attention ; la plupart des propriétés \inserted{de $W$} que
nous allons passer en revue ici, s'exprimeront par \deleted{\unsure{des}} une inclusion de la forme
\[
\mathcal F\subset\mathrm{Fib}_W\,,
\]
où $\mathcal F$ est une des sept autres classes du
diagramme (4.2.1) ; à l'exception \deleted{\unsure{tout}}
au plus du \inserted{cas $\mathcal F=$} $\mathrm{Fib}_{W_0}$, i.e. de l'inclusion
$\mathrm{Fib}_{W_0}\subset\mathrm{Fib}_W$. Comme propriété sur $W$, cette
inclusion paraît très faible et \unsure{pour|par} cela peu
utile à grand chose, sauf bien sûr dans
le cas où $\mathrm{Fib}_{W_0}\longrightarrow\mathrm{Fib}f_{W_0}$ est une égalité,
de \deleted{qui} sorte que l'inclusion \unsure{envisagée}
équivaut à $\mathrm{Fib}f_{W_0}\subset\mathrm{Fib}_W$ (et implique
$\deleted{\mathrm{Perf}_{W_0}\subset\unsure{\mathrm F|\mathbb F}}$ que $\mathrm{Perf}_{W_0}$ et $\mathcal F_{W_0}$ sont eux aussi
contenus dans $\mathrm{Fib}_W$). C'est là par contre,
comme nous allons voir, une propriété assez
forte et fort utile de $W$, impliquant notamment le fait que $(\mathrm{Cat},W,\mathrm{Fib}_W)$ est une
catégorie à fibrations.
% archive sheet number: 306
\scanpage{308}\mspage{32}
Aux fins de \unsure{référence|références} \unsure{ultérieure|ultérieures}, je vais
formuler maintenant les divers\deleted{\unsure{es}} aspects principaux
de l'\uline{axiome des} \inserted{$W$-}\uline{fibrations} $\mathrm{Loc}(7)$. La formulation
initiale à laquelle j'étais parvenu
(cf $\overline{\underline{\mathrm I}}$, $\overline{\underline{\mathrm{VI}}}$, $\overline{\underline{\mathrm{XII}}}$) \deleted{\unsure{du}} utilisait la notion
\unsure{que} j'appelle à présent ``$W$-fibration \uline{faible}'' \deleted{\unsure{,}}
(\inserted{l'}ex-$W$-fibration ; désormais détrônée !), et
sous sa forme la plus forte s'écrivait comme
l'inclusion
\[
\mathrm{Fib}f_W\subset\mathrm{Fib}_W\;.
\]
Cette notion de $W$-fibration faible, tout comme
celle de morphisme $W$-parfait, si \unsure{utile|utiles}
\unsure{soient-elles|soient-elle}, m'apparaît cependant un
peu technique, et je préfère donner
une formulation qui ne fasse appel
qu'à des notions plus familières et
d'apparence plus élémentaire. C'est pourquoi
je \deleted{\unsure{préfère me}} vais remplacer $\mathrm{Fib}f_W$ par
la sous-classe $\mathcal F_W\unsure{,}=\deleted{\unsure{\mathrm{Fib}}}\mathrm{Catfib}^{W}_{\deleted{\unsure{W}}}$, \deleted{\unsure{Relativement}} \inserted{ce qui}
(on le verra) \inserted{\unsure{revient}} \deleted{\unsure{naturellement dans l'optique}} en fait au même
% archive sheet number: 307
\scanpage{309}\mspage{33}
Cette classe $\mathcal F_W$ est extrêmement naturelle, et
s'impose même dans \deleted{\unsure{l'optique}} \inserted{\unsure{le} yoga} des dérivateurs.
En effet, dans la construction de
$\mathrm{HOT}_W(S)$ pour $S$ dans $\mathrm{Cat}$, on passe
à une catégorie de fractions de la catégorie
\[
\deleted{\unsure{\mathrm{Cat}}}\qquad\mathrm{Cat}(S)=\underline{\mathrm{Hom}}\,(S,\mathrm{Cat})\,,
\]
par un localiseur $W(S)$ (que nous
avions généralement noté $W^{g}_S$ \inserted{ou $W^{\unsure{u|i}}_S$} \ldots). Cette
catégorie \deleted{est} $\mathrm{Cat}(S)$ est d'autre part
isomorphe à la catégorie
\[
\deleted{\unsure{\mathrm{Cat}}}\mathrm{Fibsc}\ S
\]
des catégories fibrées et scindées sur $S$, avec
comme flèches les $S$-foncteurs qui respect\deleted{\unsure{ant}}ent
les scindages (et en particulier cartésiens)
(\unsure{--} cf VI, et notamment le diagramme p. 175).
On a vu (loc. cit) que \deleted{\unsure{cette}} la catégorie
localisée $\mathrm{HOT}_W(S)$ est équivalente à la
catégorie localisée de la catégorie très voisine
\[
\begin{aligned}
\mathrm{Fib}\, S\quad &\text{(que je préférerais maintenant noter }\mathrm{Catfib}\,S\ \text{ou}\\
&\mathrm{Catfib}_S)\,,
\end{aligned}
\]
formée des catégories \uline{fibrées} sur $S$, avec comme
% archive sheet number: 308
\scanpage{310}\mspage{34}
flèches les seuls foncteurs cartésiens. (On fera
attention cependant que le foncteur canonique
$\mathrm{Fibsc}\,S\mathrel{\unsure{\text{---}|\longrightarrow}}\mathrm{Fib}\,S$ est fidèle mais non pleinement ; néanmoins il induit une équivalence sur les catégories localisées, \deleted{\unsure{)}} en
définissant de façon évidente, par \unsure{la|les} $W$-\unsure{équivalence|équivalences} fibre par fibre, i.e comme $W_S^u$,
le localiseur \inserted{pertinent} sur $\mathrm{Fib}\,S$.) Donc pour cette
raison de principe, et \deleted{\unsure{d'a}}d'ailleurs tout au
long du développement de la théorie
de $\mathrm{HOT}_W$, la catégorie $\mathrm{Fib}\,S$ et par là,
les foncteurs fibrants, ont joué un rôle
essentiel. La \deleted{catégorie} \inserted{classe de flèches} que nous notons ici
$\deleted{\unsure{\mathcal F}}\mathcal F_W$ ou $\mathrm{Catfib}^{W}_{\deleted{\unsure{W}}}$ (on pourrait aussi la
noter $\mathrm{Fib}^{W}_{\deleted{\unsure{W}}}$, \deleted{\unsure{si}} ce n'était pour les risques
de confusion avec $\mathrm{Fib}_W$) \deleted{\unsure{est}} s'\deleted{\unsure{ex}}\unsure{interprète|explicite}
alors très naturellement à partir de là.

Une flèche $f:X\longrightarrow S$ est dans $\mathcal F_W$ sss
elle est \quad a) fibrante, et par là définit
de façon particulièrement directe un objet
de $\mathrm{HOT}_W(S)$, puisque à $S$-équivalence près, $X/S$
provient d'un \deleted{\unsure{élément}} \inserted{objet} de $\mathrm{Fibsc}\,S$ ou encore de $\mathrm{Cat}(S)$),
% archive sheet number: 309
\scanpage{311}\mspage{35}
et \quad b) l'objet dans $\mathrm{H}\deleted{\unsure{\mathrm o}}\mathrm{OT}_W(S)$ défini par $f$ est
dans $\mathrm{HOT}^{lc}_W(S)$, i.e. dans l'image essentielle
de la sous-catégorie str. pleine $\mathrm{Fibsc}^{W}_{\unsure{\circ|0}}\,S$ de
$\mathrm{Fibsc}\,S$ (par le foncteur canonique \deleted{\unsure{à}}
$\mathrm{Fibsc}\longrightarrow\mathrm{HOT}_W(S)$), ou ce qui revient au
$=$ (du moins si $W$ est fortement saturé)
que le foncteur
\[
\begin{aligned}
S^{\circ}&\longrightarrow\mathrm{Hot}_W\\
s&\longmapsto\mathrm{hot}_W(X_s)
\end{aligned}
\]
soit \deleted{\unsure{isomorphe à}} un foncteur \inserted{``localement} constant''. Ainsi,
on \unsure{considèrera|considérera} les $f\in\mathcal F_W$ comme les
flèches les plus économiques et les plus
naturelles dans $\mathrm{Cat}$, qui définissent des
$W$-types d'homotopie relatifs localement constants.

Notons aussi que toutes les \deleted{catégories} \inserted{classes} de
flèches (4.2.1) sont stables par chgt de base.
(\deleted{\unsure{et}} et elles contiennent les iso\unsure{)}. Les catégories
$\mathrm{Fib}_W$, $\mathrm{Perf}_W$ et leurs variantes $W_0$ sont aussi
stables par composition, par contre ce n'est
pas le cas pour $\mathrm{Fib}f_W$ et $\mathcal F_W$ et leurs
variantes $W_0$, sauf axiomes particuliers sur
$W$ \unsure{et|ou} sur $W_0$,) Il en résulte aussitôt
% archive sheet number: 310
\scanpage{312}\mspage{36}
que si $\mathcal F$ est une de ces \deleted{catégories} \inserted{classes}, la
condition
\[
\mathcal F\subset\deleted{\unsure{\mathcal F}}\mathrm{Fib}_W
\]
équivaut à celle-ci :

Pour \deleted{\unsure{tout}} diagramme cartésien
\diagram{\begin{tikzcd}[ampersand replacement=\&,column sep=normal,row sep=small]
\& X\arrow[d,dash,"f"'] \& X'\arrow[l,"q"']\arrow[d,"f'"]\\
(4.2.2) \& S \& S'\arrow[l,"p"]\text{,}
\end{tikzcd}}
avec $f\in\mathcal F$ et $p\in W$, on \unsure{á|a} $q\in W$. (NB, on
a aussi $f'\in\mathcal F$, par \inserted{la} stabilité rappelée plus haut.)

4.\deleted{\unsure{2}}3. Voici donc \inserted{(cf plus bas)} l'axiome pertinent, $\mathrm{loc}(7)$
\deleted{\unsure{avec}} en trois aspects $a,b,c$\textsuperscript{*}. On
% slanted
\margin{j'y ai rajouté finalement une quatrième, $\deleted{\unsure{\mathrm F}},\unsure{\varphi|q}$ (comme ``\unsure{fibration|fibrations}'')}
a vu (et on verra) que l'on a les
implications
\[
a\Longrightarrow b\Longrightarrow c.\,[\Longrightarrow\unsure{\mathrm F|\varphi}].
\]
Quand on référera à $\mathrm{Loc}(7)$ sans autre précision,
il sera entendu qu'il s'agit de la
forme la plus forte $\mathrm{loc}(7a)$, qui implique
les deux autres. Il n'en faut pas
\deleted{\unsure{tt}} \unsure{tellement, à l'aire|tellement, à vrai dire}. Mais
\unsure{je rappelle|rappelons} que la variante \inserted{faible} \deleted{\unsure{Loc}} $(7b)$ est
\unsure{suffisante} pour faire la théorie des foncteurs
$\deleted{\unsure{f_!}} f_!$ pour les catégories $\mathrm{HOT}_W(S)$\textsuperscript{*} (cf $\overline{\underline{\mathrm{VI}}}$).
% slanted
\margin{Ici ma mémoire me joue un tour -- on n'a nullement besoin pour cela de l'\unsure{axiome} déjà très fort $\mathrm{loc}(7b)$, mais seulement de $\mathrm{Loc}(3\ \mathrm{bis})$, dont il \unsure{semble bien} qu'il soit vérifié pour tous les localiseurs fond. ``raisonnables''. \deleted{\illegible{1}}}
% archive sheet number: 311
\scanpage{313}\mspage{37}
C'est d'ailleurs une chose surprenante,
puisque via le yoga général des dérivateurs
associés aux catégories de modèles, on s'attendrait plutôt que pour avoir des adj. à g.
$f_!$ des foncteurs $f^*$ d'image inverse, il faut
faire sur $W$ des hypothèses du type
\uline{$W$-cofibration} \deleted{\unsure{(p. ex.}} (essentiellement, $\mathrm{loc}(6a)$).
Il me semble peu probable que \deleted{\unsure{cet}} $\mathrm{Loc}(7b)$,
ni même $\mathrm{Loc}(7)$, implique $\mathrm{Loc}(6a)$, ni
même sa variante \inserted{plus} faible $\mathrm{Loc}(6'a)$ (axiome
du carré cocartésien, ou de Mayer-Vietoris,
forme faible). \deleted{\unsure{En l'absence de}} Si d'autre
part ces axiomes ne sont pas \uline{vérifiés}, alors
que $\mathrm{Loc}(7)$ l'est, il y a des chances que
cela donnera \inserted{un} \deleted{\unsure{l'}}exemple \deleted{\unsure{d'un}} où il
y a un formalisme \inserted{des} $f_!$, sans que
l'axiome d'exactitude à gauche soit
vérifié.

Quant à la forme $\mathrm{Loc}(7c)$, plus faible
encore, elle suffit pourtant (comme nous
le verrons plus bas) à assurer que
$\mathrm{Cat}$ admet suffisamment de $W$-fibrations
% archive sheet number: 312
\scanpage{314}\mspage{38}
(résultat qui n'a rien d'évident, et ne pouvait
m'apparaître qu'après avoir pris connaissance
du substantiel lemme de K. \deleted{\unsure{A.}}S. Brown,
dans son beau travail où il introduit les
catégories à fibrations et à cofibrations).

Nous donnerons aussi plus bas des cas
fort simples, où la condition $\mathrm{Loc}(7b)$ \deleted{\unsure{sous}}
entraîne \inserted{déjà} $\mathrm{Loc}(7)$ ($W$ \unsure{fortmt} saturé, i.e.
$\mathrm{Loc}(1\ \mathrm{bis})$.), voire même où $\mathrm{Loc}(7c)$
entraîne $\mathrm{Loc}(7)$ (cas où $W$ satisfait à
l'axiome \deleted{\unsure{(7)}} $\mathrm{Loc}(8)$ des limites inductives).
Je signale en passant que je ne connais
pas de localiseur \deleted{\unsure{satisf}}fondamental satisfaisant
l'axiome $\mathrm{Loc}(8)$ (qui implique déjà $W=W^{\unsure{\leftrightarrows|\rightleftarrows}}$ \inserted{i.e. $\mathrm{loc}(1\ \mathrm{bis})$})
sans que $W$ soit \unsure{fortmt} saturé \deleted{\unsure{(}} (i.e. ne
\unsure{satisfasse} même $\mathrm{Loc}(1\ \unsure{\mathrm{ter}|\mathrm{bis}})$).

\uline{Loc 6} \quad a) \deleted{\unsure{Soit}} On a $\deleted{\unsure{\mathcal F}}\mathcal F_W\subset\deleted{\unsure{\mathrm{Fib}f}}\mathrm{Fib}_W$, i.e. pour
tout diagramme \unsure{cartésien} (4.2.2) dans $\mathrm{Cat}$,
avec $f\in\mathcal F_W$ (i.e. $f$ Cat-fibrant et \unsure{fblt.} $W$-fibrant),
et $p\in W$, on a $\unsure{q|f}\in W$.

b) Pour \unsure{tout} diagramme comm.
\diagram{\begin{tikzcd}[ampersand replacement=\&,column sep=small,row sep=small]
X\arrow[rr,dash,"f"]\arrow[dr,dash,"p"'] \& \& Y\arrow[dl,dash,"q"]\\
\& \unsure{\bullet|s}\, S \&
\end{tikzcd}}
% archive sheet number: 313
\scanpage{315}\mspage{39}
$p,q\in\mathcal F_W$ et $f\in W$, on a $f\in W^{\unsure{u|v}}$ (ou,
ce qui revient au même, \deleted{\unsure{cf $\overline{\underline{\mathrm{VI}}}$}} comme $X,Y$ sont
$W$-lisses sur $S$, $f_s\in W$ pour tt $s\in\mathrm{Ob}\, S$).
% slanted
\margin{\uline{NB} La \unsure{dém.} de $\mathrm B(\mathcal F_W)\Longrightarrow\mathrm B(\mathrm{Fib}f_W)$ montre en fait qu'il suffit de postuler b) lorsque $f$ est cartésien, et $=$ lorsque $X,Y$ sont scindées sur $S$ et $f$ compatible avec les scindages.}

c) \deleted{\unsure{Pour tte flèche}} On a
\[
\mathcal F_W\cap W\subset W^u\,,
\]
en d'autres termes, pour toute flèche
$f:X\unsure{\mathrel{\text{---}}|\longrightarrow} S$ dans $\mathrm{Cat}$ telle que $f\in\mathcal F_W$
(on pourrait appeler ces flèches des Cat-$W$-fibrations, étant entendu que \unsure{cela|elles} sont \inserted{aussi} des
Cat-fibrations, mais non \deleted{\unsure{f}}nécessairement
des $W$-fibrations), si $f\in W$, alors
$f\in W_S^u$ (ou, ce qui revient au $=$, les fibres
$X_s$ de $f$ ($s\in\mathrm{Ob}\,S$) sont $W$-asphériques).

% slanted
\margin{rajouté après coup $\rightsquigarrow$}
\inserted{F) $\mathrm{Cat}$ a suffisamment de $W$-fibrations, i.e. $\mathrm{Cat}$
munie de $(W,\mathrm{Fib}_W)$ est une cat. à cofibrations (propres, et
$=$ de \deleted{\unsure{Br}}\unsure{Brown}, i.e. $X\longrightarrow e$ est dans $\mathrm{Fib}_W$ pour tt $X$ dans $\mathrm{Cat}$).}

\unsure{Il} est clair que $\unsure{b|c}$ est \unsure{un} cas particulier
de $\unsure{c|b}$ obtenu en prenant $Y=S$, $q=\mathrm{id}_S$,
donc $b\Longrightarrow c$. L'implication $a\Longrightarrow b$ est
moins évidente, elle est établie d'abord dans
$\overline{\underline{\mathrm{XII}}}$, et nous y reviendrons ci-dessous \inserted{(4.4.2 ci-dessous)}.
Le reste du présent \S\ est consacré à \deleted{\unsure{donner}} \inserted{étudier}
diverses variantes des conditions a) b) c)
précédentes, et les implications \deleted{entre} celles-ci.
% archive sheet number: 314
\scanpage{316}\mspage{40}
4.4. Soit, plus généralement
\[
\mathcal F\subset\mathrm{Fl}(\mathrm{Cat})
\]
un ens. de flèches dans $\mathrm{Cat}$, en pensant
plus particulièrement au cas où $\mathcal F$ est un
des sept \deleted{catégories} \inserted{ens.} du diagramme
(4.2.1.) (p.~30) \deleted{\unsure{distinctes de|distincts de}} \inserted{autres que} $\mathrm{Fib}_W$.
\deleted{Supposons que $\mathcal F$ contienne les iso et
soit stable par changement de base.} Nous
pouvons formuler \deleted{des} \inserted{trois} propriétés similaires
pour $\mathcal F$ à celles formulées dans $\mathrm{Loc}(7,abc)$
dans le cas de $\mathcal F_W$, notons ces
propriétés
\[
A(\mathcal F),\ B(\mathcal F),\ C(\mathcal F),
\]
% slanted
\margin{\uline{NB} Ces conditions ont un sens si $\mathcal F\subset\mathrm{Fl}(M)$, si $(M,W)$ est une catégorie de modèles quelconque}
donc
\[
\left\{\begin{aligned}
A(\mathcal F):&\quad\mathcal F\subset\mathrm{Fib}_W\\
B(\mathcal F):&\quad\mathrm{Fl}(\mathcal F_S)\cap W_{\unsure{/S|S}}\subset W_S^u\\
&\quad\text{(où $\mathcal F_S$ désigne la sous-catégorie pleine de}\\
&\quad\text{$\mathrm{Cat}_{/S}$ formée des $f:X\text{---}S$ avec $f\in\mathcal F$)}\\
C(\mathcal F):&\quad\mathcal F\cap W\subset W^u.
\end{aligned}\right.
\tag*{(4.4.1)}
\]
Il est clair qu'on a $B(\mathcal F)\Longrightarrow C(\mathcal F)$, et je
vais \deleted{\unsure{montrer}} % Broken underline below both parts of the cancelled phrase.
\uline{\deleted{donner une condition sur $\mathcal F$ qui
implique}} qu'on a aussi $A(\mathcal F)\Longrightarrow B(\mathcal F)$,
% archive sheet number: 315
\scanpage{317}\mspage{41}
[Nous l'appliquerons dans les cas suivants
\[
\mathcal F=\mathcal F_{W^{\deleted{\unsure{k|*}}}},\mathrm{Fib}f_{W^{\deleted{\unsure{k|*}}}},\mathrm{Perf}_W,
\deleted{\unsure{\mathrm{Perf}_{W_0}}}\mathcal F_{W_0^{\deleted{\unsure{k|*}}}},
\mathrm{Fib}f_{W_0^{\unsure{k|*}}},\mathrm{Perf}_{W_0}.
\tag*{(4.4.1)}
\]
Je signale d'ailleurs que l'on verra plus bas \inserted{$\underbrace{\text{(cor. }\unsure{4.4.10}\text{)}}$} que
\[
 A(\mathcal F_W)\Longleftrightarrow A(\mathrm{Fib}f_W),\quad
 A(\mathcal F_{W_0^{\unsure{k|*}}})\Longleftrightarrow A(\mathrm{Fib}f_{W_0^{\unsure{k|\alpha|*}}}),
 \tag*{(4.4.3)}
\]
et de même pour les variantes $B,C$, de
sorte que pour l'application que nous voulons
étudier, il n'y a pas lieu de distinguer les
$\mathcal F_W$ et $\mathrm{Fib}f_W$, les $\mathcal F_{W_0}$ et les $\mathrm{Fib}f_{W_0}$.
D'ailleurs, il suffira de regarder les \deleted{\unsure{catég.}}
ensembles $\mathcal F_{W_0},\mathrm{Fib}f_{W_0},\mathrm{Perf}_{W_0}$, puisque \unsure{le|les}
cas \deleted{des} \deleted{\unsure{catég.}} \unsure{relatif|relatifs} à $W$ en sont un cas
particulier (en prenant $W_0=W$).]
% slanted
\margin{Cela fera donc a priori $3\times6=18$ conditions différentes, mais à cause de (4.4.3) il n'y en a (à \unsure{équiv.|isom.|ignor.} près) que $4\times3=12$, et en fait deux parmi ces quatre triplets de conditions se réduisent à une seule (pour $\mathcal F=\mathcal F_W$ et $\mathcal F=\mathcal F_{W_0}$), de sorte qu'il ne reste au total que \uline{huit} conditions différentes, que nous \unsure{résumons|désignons} dans le diagramme récapitulatif (\unsure{4.5.1|4.5.1.1}) de la page \unsure{48|40}. Il y en a \deleted{\unsure{en}} seulement \uline{quatre} $A_{W_0},B_{W_0},C_{W_0},D_{W_0}$ \unsure{comme|à donner comme} fonctions de $W_0$, les quatre autres $A,B,C,D$ s'obtiennent en \unsure{en} spécialisant au cas où $W_0=W$.}

\deleted{Cela résultera évidemment du}

\deleted{\uline{Lemme} 4.4.1 Soit $f:X\text{---}Y$ un $S$-morphisme dans $\mathrm{Cat}$, avec $p:X\text{---}S$ et $q:Y\longrightarrow S$ des $W$-fibrations. Si $f\in W$, donc $f\in W^u$. Cela résulte en effet formellement du fait que $\mathrm{Cat}$ a assez de $W$-fibrations}
% slanted
\margin{Vaut dans une catégorie de modèles quelconque.}

\uline{Lemme} 4.4.1. Supposons que $\mathcal F$ satisfasse
la condition suivante : toute flèche $f$ de
$\mathrm{Cat}$ se factorise en $f=pi$, avec $i\in W$, $f\in\mathcal F$.
Alors la condition $A(\mathcal F)$ sur $\mathcal F$ implique
que dans $\mathrm{Cat}$ il y a ``assez de $W$-fibrations'',
et par suite (par un argument formel bien
% archive sheet number: 316
